\begin{abstract}
Knowledge graph repair has two linked tasks: get useful rules and fix graph
errors within a budget. We train a Double-DQN policy to choose between them.
The test has four graph actions and four rule actions in a fixed rule library.
Two text prompts also generate rules. One removes short text spans; the other
adds new sentences. We test the generated rules in saved-output repair runs.
An action test shows that a penalty based on error counts blocks useful
relation repair. The repair can create temporary duplicate edges.
We compare count, rate, and zero penalties using forty neural models and
ten one-step fitted baselines. On thirty new corruptions of the same base
graph, the rate penalty raises Double-DQN triple F1 from 94.67\% to 99.49\%.
It restores 94.00\% of the original invalid relations.
A simple acquire-then-deficit rule reaches 99.91\% F1.
The 9,456 saved model responses contain different candidates from the two
prompts. With equal call budgets, augmentation yields more constraints.
We compile 18,143 exact type patterns. Their union detects 10 of 64 defects
in RuleTest-94. The repair loop removes these ten records and keeps all
thirty clean records. Further document tests count recovered facts and
correct facts lost. The results show how reward design changes repair
and how rule coverage shapes the choice of later actions.
\end{abstract}
